\section{ATFM problem optimization}
\begin{frame}{}

    Terminal Manoeuvering Area aircraft conflict-free routing and landing sequencing to solve the multi-objective ATFM\footnote{Air Traffic Flow Management} large scale problem at TMA level.
    \vfill
    \begin{flushleft}\scriptsize
        \colorea{colorC}{eric.lucero@alumnos.upm.es\\
            Started: 2025}
    \end{flushleft}
\end{frame}
\begin{frame}{TMA problem}
    \begin{itemize}
        \item<+-> \bfb{Problem/Question}\\
              TMA congestion and conflict are uprising due to the increment of passengers/cargo flight numbers. That growing demand has so many
              consequences:\\
              \begin{itemize}[itemsep=1mm]\small
                  \item Safety issues
                  \item Delays in arrival/departure
                  \item Congestion in TMA
                  \item ATCO\footnote{Air Traffic Controller} workload
              \end{itemize}
    \end{itemize}
    \small
    \begin{questionBox}
        {Question}{
            % Model the TMA for a given number of arriving/departure time that simulates the real  .... of an airport and optimize the delays and number of conflicts using a new metaheuristic algorithm. \\
            How can air traffic in the Terminal Manoeuvering Area (TMA) be modeled and optimized to reduce delays and minimize conflicts under high air traffic demand conditions using a novel multi-objective metaheuristic algorithm?}
    \end{questionBox}
\end{frame}
\begin{frame}
    \begin{itemize}
        \item<+-> \bfb{Observation/Investigation}\\
              The increasing air traffic demand has led to rising congestion and conflicts in the TMA, resulting in significant operational challenges. These challenges include safety concerns, delays in arrivals and departures, increased congestion, and higher workload for air traffic controllers (ATCOs). Traditional optimization methods struggle to efficiently balance these conflicting objectives. This study investigates how a novel multi-objective metaheuristic algorithm can enhance traffic flow management, improve scheduling efficiency, and reduce overall system delays and conflicts in real-world TMA scenarios.
    \end{itemize}

\end{frame}
\begin{frame}{TMA problem}
    \begin{figure}
        \centering
        \includegraphics[width=0.95\textwidth]{2024_zhang_A_Study}
        \caption*{\tiny\citeauthor{2024_Yi_Zhang}; \citetitle{2024_Yi_Zhang}, \citeyear{2024_Yi_Zhang}}
    \end{figure}
\end{frame}


\begin{frame}
    \begin{itemize}
        \item<+-> \bfb{Hypotheses}\\
              \begin{questionBox}
                  {Hypotheses}{
                      A \red{novel} metaheuristic based on \red{...} will enable the reduction of delays and conflicts in the \red{Terminal Manoeuvring Area (TMA)}
                      by at least \red{5 \%} compared to traditional methods, while optimizing safety, conflict resolution and sector congestion in a multi-objective framework.
                  }
              \end{questionBox}
        \item<+-> \bfb{Experiments/Data}\\
              Test the new metaheuristic against existing ones trough benchmark test, and validate the results of the TMA problem optimization with real flight data.
    \end{itemize}

\end{frame}